% Preamble injected before each .tex file so pandoc can handle
% Python-documentation-specific LaTeX macros.

% Inline code / identifier styles - all become \texttt
\newcommand{\code}[1]{\texttt{#1}}
\newcommand{\module}[1]{\texttt{#1}}
\newcommand{\function}[1]{\texttt{#1}}
\newcommand{\class}[1]{\texttt{#1}}
\newcommand{\exception}[1]{\texttt{#1}}
\newcommand{\data}[1]{\texttt{#1}}
\newcommand{\file}[1]{\texttt{#1}}
\newcommand{\kbd}[1]{\texttt{#1}}
\newcommand{\samp}[1]{\texttt{#1}}
\newcommand{\var}[1]{\emph{#1}}
\newcommand{\env}[1]{\texttt{#1}}
\newcommand{\refmodule}[2][]{\texttt{#2}}
\newcommand{\member}[1]{\texttt{#1}}
\newcommand{\method}[1]{\texttt{#1}}
\newcommand{\attr}[1]{\texttt{#1}}
\newcommand{\ctype}[1]{\texttt{#1}}
\newcommand{\cfunc}[1]{\texttt{#1}}
\newcommand{\cmacro}[1]{\texttt{#1}}
\newcommand{\sectcode}[1]{\texttt{#1}}
\newcommand{\programopt}[1]{\texttt{#1}}
\newcommand{\longprogramopt}[1]{\texttt{-{}-#1}}
\newcommand{\pep}[1]{PEP #1}
\newcommand{\rfc}[1]{RFC #1}
\newcommand{\newsgroup}[1]{\texttt{#1}}
\newcommand{\character}[1]{\texttt{#1}}
\newcommand{\regexp}[1]{\texttt{#1}}
\newcommand{\token}[1]{\texttt{#1}}
\newcommand{\production}[2]{\texttt{#1} ::= #2}
\newcommand{\productioncont}[1]{\qquad #1}
\newcommand{\grammartoken}[1]{\texttt{#1}}

% Version notes
\newcommand{\shortversion}{}
\newcommand{\version}{}
\newcommand{\release}[1]{}
\newcommand{\versionadded}[2][]{\emph{New in version #2.}}
\newcommand{\versionchanged}[2][]{\emph{Changed in version #2.}}
\newcommand{\deprecated}[2]{\emph{Deprecated since version #1: #2}}

% Text style macros
\newcommand{\dfn}[1]{\emph{#1}}
\newcommand{\strong}[1]{\textbf{#1}}
\newcommand{\emptyline}{}

% Named constants
\newcommand{\UNIX}{Unix}
\newcommand{\Cpp}{C++}
\newcommand{\C}{C}
\newcommand{\ABC}{ABC}
\newcommand{\EOF}{EOF}
\newcommand{\NULL}{NULL}
\newcommand{\POSIX}{POSIX}
\newcommand{\ASCII}{ASCII}
\newcommand{\Python}{Python}
\newcommand{\email}[1]{\texttt{#1}}
\newcommand{\url}[1]{\texttt{#1}}
\newcommand{\citetitle}[2][]{#2}

% Block-code wrappers used in early docs (1.0.1-1.4)
\newcommand{\bcode}{}
\newcommand{\ecode}{}

% Module declaration - discard
\newcommand{\declaremodule}[3][]{}
\newcommand{\modulesynopsis}[1]{\emph{#1}}
\newcommand{\platform}[1]{}

% Index commands - discard (no index in web docs)
\newcommand{\index}[1]{}
\newcommand{\indexii}[2]{}
\newcommand{\indexiii}[3]{}
\newcommand{\indexiv}[4]{}
\newcommand{\kwindex}[1]{}
\newcommand{\stindex}[1]{}
\newcommand{\exindex}[1]{}
\newcommand{\obindex}[1]{}
\newcommand{\bifuncindex}[1]{}
\newcommand{\modindex}[1]{}
\newcommand{\bimodindex}[1]{}
\newcommand{\stmodindex}[1]{}
\newcommand{\exmodindex}[1]{}
\newcommand{\refmodindex}[1]{}
\newcommand{\refbimodindex}[1]{}
\newcommand{\refexmodindex}[1]{}
\newcommand{\refstmodindex}[1]{}
\newcommand{\setindexsubitem}[1]{}
\newcommand{\withsubitem}[2]{#2}
\newcommand{\ttindex}[1]{}
\newcommand{\opindex}[1]{}

% Misc layout / navigation commands - discard
\newcommand{\linebreak}{}
\newcommand{\sectionauthor}[2]{}
\newcommand{\moduleauthor}[2]{}
\newcommand{\maintainer}[2]{}
\newcommand{\makemodindex}{}
\newcommand{\localmoduletable}{}
\newcommand{\nodename}[1]{}
\newcommand{\filenq}[1]{\texttt{#1}}

% See-also commands
\newcommand{\seetext}[1]{#1}
\newcommand{\seealsotext}[1]{#1}
\newcommand{\seetitle}[3][]{#2}
\newcommand{\seelink}[3]{#3}
\newcommand{\seemodule}[2][]{\texttt{#1} --- #2}
\newcommand{\seeurl}[2]{\texttt{#1}: #2}

% ---------------------------------------------------------------------------
% Description environments for Python APIs
% Each renders the signature as a level-4 heading (####)
% ---------------------------------------------------------------------------

% Python function: \begin{funcdesc}{name}{args}
\newenvironment{funcdesc}[2]{%
  \subsubsection*{\texttt{#1}(#2)}%
}{}

% Python function (no index): same structure
\newenvironment{funcdescni}[2]{%
  \subsubsection*{\texttt{#1}(#2)}%
}{}

% Python method: \begin{methoddesc}[class]{name}{args}
\newenvironment{methoddesc}[3][]{%
  \subsubsection*{\texttt{#2}(#3)}%
}{}

\newenvironment{methoddescni}[3][]{%
  \subsubsection*{\texttt{#2}(#3)}%
}{}

% Python class: \begin{classdesc}{name}{args}
\newenvironment{classdesc}[2]{%
  \subsubsection*{class \texttt{#1}(#2)}%
}{}

\newenvironment{classdescni}[2]{%
  \subsubsection*{class \texttt{#1}(#2)}%
}{}

% Python exception: \begin{excdesc}{name}
\newenvironment{excdesc}[1]{%
  \subsubsection*{exception \texttt{#1}}%
}{}

% Python exception class: \begin{excclassdesc}{name}{args}
\newenvironment{excclassdesc}[2]{%
  \subsubsection*{exception \texttt{#1}(#2)}%
}{}

% Python data/attribute: \begin{datadesc}{name}
\newenvironment{datadesc}[1]{%
  \subsubsection*{\texttt{#1}}%
}{}

\newenvironment{datadescni}[1]{%
  \subsubsection*{\texttt{#1}}%
}{}

% Python member: \begin{memberdesc}[class]{name}
\newenvironment{memberdesc}[2][]{%
  \subsubsection*{\texttt{#2}}%
}{}

\newenvironment{memberdescni}[2][]{%
  \subsubsection*{\texttt{#2}}%
}{}

% Python opcode: \begin{opcodedesc}{name}{args}
\newenvironment{opcodedesc}[2]{%
  \subsubsection*{\texttt{#1}\ #2}%
}{}

% C function: \begin{cfuncdesc}{returntype}{name}{args}
\newenvironment{cfuncdesc}[3]{%
  \subsubsection*{\texttt{#2}(#3)}%
}{}

% C type: \begin{ctypedesc}{name}
\newenvironment{ctypedesc}[1]{%
  \subsubsection*{\texttt{#1}}%
}{}

% C variable: \begin{cvardesc}{type}{name}
\newenvironment{cvardesc}[2]{%
  \subsubsection*{\texttt{#2}}%
}{}

% C macro: \begin{csimplemacrodesc}{name}
\newenvironment{csimplemacrodesc}[1]{%
  \subsubsection*{\texttt{#1}}%
}{}

\newenvironment{csimplemacro}[1]{%
  \subsubsection*{\texttt{#1}}%
}{}

% C macro with args: \begin{cmacrodesc}{name}{args}
\newenvironment{cmacrodesc}[2]{%
  \subsubsection*{\texttt{#1}(#2)}%
}{}

% Env var: \begin{envdesc}{name}
\newenvironment{envdesc}[1]{%
  \subsubsection*{\texttt{#1}}%
}{}

% Macro desc: \begin{macrodesc}{name}{args}
\newenvironment{macrodesc}[2]{%
  \subsubsection*{\texttt{#1}(#2)}%
}{}

% ---------------------------------------------------------------------------
% Container environments
% ---------------------------------------------------------------------------

% See-also block
\newenvironment{seealso}{%
  \paragraph*{See also:}%
}{}

% definitions / fulllineitems act like description lists
\newenvironment{definitions}{\begin{description}}{\end{description}}
\newenvironment{fulllineitems}{\begin{description}}{\end{description}}

% synopsistable -- module synopsis table, render as simple tabular
\newenvironment{synopsistable}{\begin{tabular}{ll}}{\end{tabular}}

% longtableii / longtableiii: Python 1.5+ long tables with headers
% Args: {colspec}{headfmt}{col1head} etc. -- just render as tabular
\newenvironment{longtableii}[3]{\begin{tabular}{#1}}{\end{tabular}}
\newenvironment{longtableiii}[4]{\begin{tabular}{#1}}{\end{tabular}}
\newenvironment{longtableiv}[5]{\begin{tabular}{#1}}{\end{tabular}}

% lineheader inside longtableii (used as header row)
\newcommand{\lineheader}[1]{\textbf{#1}}
\newcommand{\lineitemi}[2]{#1 & #2 \\}
\newcommand{\lineitemii}[3]{#1 & #2 & #3 \\}
\newcommand{\lineitemiii}[4]{#1 & #2 & #3 & #4 \\}

% sloppypar is standard LaTeX; define just in case
\newenvironment{sloppypar}{}{}
